\section{Verification method}\label{sec:verification}

A sampler that is fast but subtly biased is worse than a slow one, because its errors propagate silently into
threshold estimates. The kernel is therefore checked at every level of its implementation against a single software
model, and the model is checked against Stim.

\paragraph{Software model.} The model interprets the compiled program in Python with one boolean frame per qubit and
shot. Its gate, measurement, noise and fold semantics are those of the model of the first kernel, which was validated
against Stim and against single fault injection from Stim's detector error model (Section~\ref{sec:first}); only the
random stream of Section~\ref{sec:noise} is new. Given the program, the seed, the node key and the batch range, the model
produces the exact output words the kernel must produce.

\paragraph{Bit exact levels.} The same comparison is applied at four levels, each catching a different class of
error:
\begin{enumerate}
\item C simulation of the kernel source compiled with g++, for four circuits (repetition code, surface code memory in
the Z and the X basis), two seeds and two batch offsets, one of them near the top of the 32 bit batch counter range. Noiseless circuits must give all zero outputs.
\item The negative control of Section~\ref{sec:frame}, which shows that the C simulation would detect a schedule
violation.
\item RTL cosimulation of the synthesised kernel with Vitis HLS, which runs the generated hardware description,
including the pipelined frame stage and the concurrent dataflow stages: the repetition code, in which \NumRepReusesAtD{} of the
\NumRepReuses{} reuses of a qubit are at the minimum schedule distance, and \NumCosimBatchesA{} and \NumCosimBatchesB{} batches of the
distance 3 surface code, which uses every gate and noise type of the benchmarks.
\item Every kernel instance of every card on hardware, for the three benchmark circuits, with the node key of that
card and batch offsets in two launches.
\end{enumerate}

\paragraph{Statistical levels.} Bit exactness against the model establishes that the hardware computes what the model
computes. Agreement with Stim then establishes that the model samples the right distribution. Two tests are run on
hardware samples. Tier 4 compares the firing rate of every detector and observable with Stim's sampler over
\NumTierFourShots{} shots per side for the three benchmark circuits, using a two proportion $z$ score per output and
requiring $|z|<5$ for all of them. Tier 5 decodes hardware samples and Stim samples of rotated surface code memory
experiments with distance $d=3,5,7$ ($d$ rounds) and physical error rate $p=0.001,0.003,0.01$ with the same
PyMatching~\cite{higgott2022pymatching,higgott2025sparse} matcher, built from Stim's detector error model, and requires
the logical error rates to agree with $|z|<5$. Distance 7 could not be run on the first kernel, which lacked the
qubit renumbering.
